\documentclass[11pt]{article}
\usepackage{amsmath,amssymb,amsthm,hyperref}
\usepackage[margin=27mm]{geometry}
\newtheorem{lemma}{Lemma}
\newcommand{\Q}{\mathbb Q}
\title{Square arguments at odd scalar primes}
\author{}\date{17 November 2026}
\begin{document}\maketitle

\section{A local simplification}

Fix distinct nonzero rational contact values $c_1,\ldots,c_k$ and
a nonzero rational pole $b$ disjoint from them. Put
\[
 R(t)=\frac{\prod_{i=1}^k(t-c_i)}{(t-b)^k}.
\]
Let $S$ contain the primes dividing the numerators or denominators
of the contacts and pole, and their differences.

\begin{lemma}
Let $p\notin S$, $\rho,q\in\Q^\times$, and suppose that
$v_p(\rho)$ is odd. Then $R(\rho q^2)$ is a $p$-adic unit.
\end{lemma}
\begin{proof}
The argument has valuation $v_p(\rho)+2v_p(q)$, which is odd
and therefore nonzero. If it is positive, every numerator and
denominator factor of $R$ is a unit. If it is negative, every factor
has the same valuation as the argument; the equal degrees cancel.
\end{proof}

Thus a single balanced contact block at argument $\rho q^2$ suffices
for the scalar-unit part of the general-source extraction proof.
For example, the two arguments
\[
 t_0=\rho q^2,\qquad
 t_1=\rho(q+T)^2
\]
both have this property at an odd scalar valuation, provided their
arguments are defined and nonzero. Their heights also give the
unconditional inequality
\[
 h(T)\le \frac{h(t_0)+h(t_1)}2+h(\rho)+\log2,
\]
since $T=(q+T)-q$ and
$2h(q)\le h(t_0)+h(\rho)$, with the same bound for $q+T$.

\section{The two remaining requirements}

This lemma does not give a reduced-variable source. The first remaining
requirement is a single detector that captures sufficient odd contact
support from both arguments. Multiplying their contact functions loses
primes where both valuations are odd. A derivative quotient recovers
equal simple collisions, but can lose a simple odd contact when the
opposite contact has a larger even valuation. In the eleven-variable
conditional source these losses are bounded using the squarefree-form
consequence of ordinary rational $abc$. That estimate is not an
unconditional input.

The second requirement is completeness. The fresh contact factors here
are quadratic forms $\rho q^2-c_i$ and
$\rho(q+T)^2-c_i$. The existing simultaneous theorem for
$N+c_iL^2$ does not directly assert simultaneous prime realization
for these quadratic norm forms with the same $q$. In particular,
the local unit lemma must not be reported as an established
twenty-three-variable construction.

The established unconditional twenty-five-variable source retains its
two separate contact blocks. Its polynomial contact-height estimate
and Lucas growth argument remain valid there. Neither is replaced by
the logarithmic radical estimate used in the eleven-variable
ordinary-$abc$ construction.

\end{document}
